\subsection{1.4 驾驶员交互}

\subsubsection{1.4.1 系统状态}

系统状态权重 \textbf{30\%}，下面两个三级指标：

这里最重要的是“车道线缺失”，占系统状态的 80\%。

\paragraph{1.4.1.1 关闭/开启/激活状态}

测试方法很直接：

\begin{itemize}
  \item 车辆启动后，在系统关闭状态下记录提示信息。
  \item 车辆启动后，开启系统但未激活，记录提示信息。
  \item 车辆启动后，系统激活，记录提示信息。
\end{itemize}

得分规则也简单：

这里考的是用户是否能知道系统当前“到底有没有在管车”。

\begin{itemize}
  \item 关闭、开启、激活不是一回事：
  \item 关闭：功能不可用或未打开。
  \item 开启：功能打开了，但条件不满足，还没真正控制车辆。
  \item 激活：系统已经在进行横向/纵向控制。
\end{itemize}

\paragraph{1.4.1.2 车道线缺失}

车道线缺失有 4 个场景，每个场景通过得 25 分，共 100 分：

得分规则：

真实道路里车道线经常磨损、断裂、被遮挡。系统不能一遇到线缺失就立刻退出，也不能退出得悄无声息。

\subsubsection{1.4.2 驾驶员监测}

驾驶员监测权重 \textbf{25\%}，下面三个三级指标：

这一项主要对应 DMS 和脱手检测。

\paragraph{1.4.2.1 驾驶员运动脱离检测}

测试场景：长直道 + 弯道组合，道路总长不小于 625 m，弯道半径不小于 250 m。试验车辆以 30 km/h 开启组合驾驶辅助功能。

测试方法：系统开启后，驾驶员先手扶方向盘但不施加握力，车辆稳定后双手脱离方向盘，同时录制方向盘、仪表和报警声音。如果车辆具备 RMF，要确保车辆在弯道内进入 RMF。

RMF 是 Risk Mitigation Function，风险减缓功能。简单说就是：驾驶员持续不响应，系统要继续控制横纵向运动，辅助车辆停到目标停车区域或进入安全状态。

得分拆成三部分：

这个规则背后有两层意思：

第一，真正脱手时必须尽快提醒。

第二，驾驶员只是轻扶方向盘时，不应该误报。 

所以它不是越敏感越好。好的系统要能区分“驾驶员真的脱离监管”和“驾驶员手在方向盘上但握力很轻”。

\paragraph{1.4.2.2 闭眼检测}

测试场景：长直道 + 弯道组合，直道长度不小于 300 m，试验车辆以 30 km/h 开启组合驾驶辅助功能。

测试方法：调整坐姿和座椅高度，让驾驶员头部完整出现在摄像头视野中。车辆稳定后，驾驶员闭眼 2 s 以上，录制仪表、驾驶员面部和报警声音。如果具备 RMF，要确保车辆在弯道内进入 RMF。

得分规则：

闭眼检测考的是疲劳/失能风险。注意它要求的是 \textbf{4 s 内报警}，不是等很久才提醒。

\paragraph{1.4.2.3 低头检测}

测试方法和闭眼类似，但动作变成：驾驶员低头看向方向盘中央 3 s 以上。

得分规则：

闭眼和低头的差别是触发时限不同：闭眼 4 s，低头 5 s。

闭眼风险更直接，所以要求更快。

\subsubsection{1.4.3 人机共驾}

人机共驾权重 \textbf{20\%}，下面两个三级指标：

它考的是控制权切换：驾驶员想接管时，系统能不能让；驾驶员完成变道/偏移后，系统能不能恢复。

测试场景：长直道，长度不小于 300 m，试验车辆以 60 km/h 激活功能。

测试方法：系统激活后，驾驶员向左或向右转动方向盘，使功能取消激活。试验车辆进入相邻车道并保持居中行驶。

\paragraph{1.4.3.1 驾驶员主动干预}

得分规则：

这项很关键。转向力矩太小就取消，说明系统太容易被误触；转向力矩太大才让权，说明系统和人抢方向。

所以它规定了一个合理窗口：\textbf{大于 2 N·m，不超过 5 N·m}。

你可以把它理解为：驾驶员明确想接管时，系统要让权；但不能轻轻碰一下就退出。

\paragraph{1.4.3.2 系统恢复能力}

得分规则：

这项考的是系统恢复逻辑。驾驶员干预后，如果车辆已经进入相邻车道并居中，系统应当能重新恢复辅助，而不是一直处于退出状态。

这里不是鼓励系统抢回控制权，而是要求：\textbf{在驾驶员完成操作、车辆进入稳定状态后，系统能恢复可用。}

\subsubsection{1.4.4 误响应}

误响应权重 \textbf{25\%}，下面三个三级指标：

这个名字很重要：不是“响应能力”，而是“误响应”。它测的是：\textbf{旁边车道有目标物时，系统不要把它误判成本车道危险而急刹。}

测试场景：至少两条车道的长直道。试验车辆驶向前方目标物，目标物在道路右侧，也就是相邻区域，不在本车路径中心。

测试参数：

得分规则：

这项非常贴近用户体验。很多辅助驾驶系统不是“不刹车”出问题，而是“乱刹车”出问题。旁边有行人、自行车、车辆时，如果系统过度保守、突然大减速，后车追尾风险和用户信任都会受影响。

所以误响应考的是目标关联能力：\textbf{目标物是不是在本车道？是否真的进入本车行驶路径？是否需要制动？}
